%%%%%%%%%%%%%%%%%%%%%%%%%%%%%%%%%%%%%%%%%%%%%%%%%%%%%%%%%%%%
% 0.7 ELEMENTARY LOGIC
% The Composition of Advanced Mathematics
%%%%%%%%%%%%%%%%%%%%%%%%%%%%%%%%%%%%%%%%%%%%%%%%%%%%%%%%%%%%

\section{Elementary Logic}
\label{sec:elementary-logic}

\prerequisite{
Familiarity with mathematical statements and notation from
\cref{sec:mathematical-language-notation}, sets from
\cref{sec:elementary-set-theory}, relations from \cref{sec:relations}, and
functions from \cref{sec:functions}.
}

\learningobjectives{
By the end of this section, the reader should be able to:
\begin{itemize}
    \item distinguish propositions from nonpropositions;
    \item identify truth values;
    \item interpret negation, conjunction, disjunction, implication, and biconditional statements;
    \item construct and interpret truth tables;
    \item distinguish converse, inverse, and contrapositive statements;
    \item interpret necessary and sufficient conditions;
    \item recognize tautologies, contradictions, and contingencies;
    \item determine logical equivalence;
    \item work with predicates and quantifiers;
    \item negate quantified statements correctly;
    \item interpret statements with multiple quantifiers;
    \item recognize common valid and invalid forms of argument;
    \item connect logical operations with set operations and proof methods.
\end{itemize}
}

%%%%%%%%%%%%%%%%%%%%%%%%%%%%%%%%%%%%%%%%%%%%%%%%%%%%%%%%%%%%
\subsection{Why Logic Matters}
\label{subsec:why-logic-matters}
%%%%%%%%%%%%%%%%%%%%%%%%%%%%%%%%%%%%%%%%%%%%%%%%%%%%%%%%%%%%

Logic studies the structure of valid reasoning.

Mathematical definitions, theorems, proofs, algorithms, and formal arguments
all rely on logical structure.

For example, when a theorem states

\[ P\Longrightarrow Q, \]

mathematics is asserting that whenever statement \(P\) is true, statement
\(Q\) must also be true.

Understanding such statements is essential before studying formal proofs.

Logic also provides a bridge between ordinary language and mathematical
reasoning.

%%%%%%%%%%%%%%%%%%%%%%%%%%%%%%%%%%%%%%%%%%%%%%%%%%%%%%%%%%%%
\subsection{Propositions}
\label{subsec:propositions}
%%%%%%%%%%%%%%%%%%%%%%%%%%%%%%%%%%%%%%%%%%%%%%%%%%%%%%%%%%%%

\begin{definitionbox}{Proposition}
A \emph{proposition} is a declarative statement that has exactly one truth
value: true or false.
\end{definitionbox}

Examples include:

\begin{itemize}
    \item \(2+3=5\);
    \item \(7<4\);
    \item every square is nonnegative;
    \item \(11\) is even.
\end{itemize}

The first and third propositions are true.

The second and fourth are false.

%%%%%%%%%%%%%%%%%%%%%%%%%%%%%%%%%%%%%%%%%%%%%%%%%%%%%%%%%%%%
\subsection{Statements That Are Not Propositions}
\label{subsec:nonpropositions}
%%%%%%%%%%%%%%%%%%%%%%%%%%%%%%%%%%%%%%%%%%%%%%%%%%%%%%%%%%%%

Questions, commands, and expressions with unspecified variables are generally
not propositions.

For example,

\begin{itemize}
    \item ``What is \(5+3\)?''
    \item ``Solve the equation.''
    \item \(x+2=7\)
\end{itemize}

are not propositions as written.

The expression

\[ x+2=7 \]

may become a proposition after a value of \(x\) is specified or after the
variable is quantified.

%%%%%%%%%%%%%%%%%%%%%%%%%%%%%%%%%%%%%%%%%%%%%%%%%%%%%%%%%%%%
\subsection{Truth Values}
\label{subsec:truth-values}
%%%%%%%%%%%%%%%%%%%%%%%%%%%%%%%%%%%%%%%%%%%%%%%%%%%%%%%%%%%%

A proposition has one of two truth values:

\[ \text{True}\qquad\text{or}\qquad\text{False}. \]

These are often abbreviated

\[ T\qquad\text{and}\qquad F. \]

In formal logic and computer science, true and false may also be represented
by other conventions such as \(1\) and \(0\), depending on context.

%%%%%%%%%%%%%%%%%%%%%%%%%%%%%%%%%%%%%%%%%%%%%%%%%%%%%%%%%%%%
\subsection{Logical Variables}
\label{subsec:logical-variables}
%%%%%%%%%%%%%%%%%%%%%%%%%%%%%%%%%%%%%%%%%%%%%%%%%%%%%%%%%%%%

Letters such as

\[ P,\qquad Q,\qquad R \]

are commonly used to represent propositions.

For example, let

\[ P:\text{ ``\(8\) is even''} \]

and

\[ Q:\text{ ``\(8>10\)''}. \]

Then \(P\) is true and \(Q\) is false.

%%%%%%%%%%%%%%%%%%%%%%%%%%%%%%%%%%%%%%%%%%%%%%%%%%%%%%%%%%%%
\subsection{Negation}
\label{subsec:negation}
%%%%%%%%%%%%%%%%%%%%%%%%%%%%%%%%%%%%%%%%%%%%%%%%%%%%%%%%%%%%

The symbol

\[ \neg \]

is called the \emph{negation symbol} and is read

\begin{center}
``not.''
\end{center}

\begin{definitionbox}{Negation}
The \emph{negation} of a proposition \(P\), written
\[ \neg P, \]
is true exactly when \(P\) is false.
\end{definitionbox}

If

\[ P:\text{ ``\(5>2\)''}, \]

then

\[ \neg P:\text{ ``\(5\) is not greater than \(2\)''}. \]

Since \(P\) is true, \(\neg P\) is false.

%%%%%%%%%%%%%%%%%%%%%%%%%%%%%%%%%%%%%%%%%%%%%%%%%%%%%%%%%%%%
\subsection{Truth Table for Negation}
\label{subsec:negation-truth-table}
%%%%%%%%%%%%%%%%%%%%%%%%%%%%%%%%%%%%%%%%%%%%%%%%%%%%%%%%%%%%

A \emph{truth table} lists all possible truth values of logical expressions.

For negation,

\[
\begin{array}{c|c}
P & \neg P\\ \hline
T & F\\
F & T
\end{array}
\]

Negation reverses the truth value.

%%%%%%%%%%%%%%%%%%%%%%%%%%%%%%%%%%%%%%%%%%%%%%%%%%%%%%%%%%%%
\subsection{Conjunction}
\label{subsec:conjunction}
%%%%%%%%%%%%%%%%%%%%%%%%%%%%%%%%%%%%%%%%%%%%%%%%%%%%%%%%%%%%

The symbol

\[ \land \]

is called the \emph{logical conjunction symbol} and is read

\begin{center}
``and.''
\end{center}

\begin{definitionbox}{Conjunction}
The conjunction
\[ P\land Q \]
is true exactly when both \(P\) and \(Q\) are true.
\end{definitionbox}

For example, the statement

\[ 5>2\land 8>3 \]

is true.

But

\[ 5>2\land 8<3 \]

is false.

%%%%%%%%%%%%%%%%%%%%%%%%%%%%%%%%%%%%%%%%%%%%%%%%%%%%%%%%%%%%
\subsection{Truth Table for Conjunction}
\label{subsec:conjunction-truth-table}
%%%%%%%%%%%%%%%%%%%%%%%%%%%%%%%%%%%%%%%%%%%%%%%%%%%%%%%%%%%%

\[
\begin{array}{c|c|c}
P & Q & P\land Q\\ \hline
T & T & T\\
T & F & F\\
F & T & F\\
F & F & F
\end{array}
\]

The conjunction is true only in the first row.

%%%%%%%%%%%%%%%%%%%%%%%%%%%%%%%%%%%%%%%%%%%%%%%%%%%%%%%%%%%%
\subsection{Disjunction}
\label{subsec:disjunction}
%%%%%%%%%%%%%%%%%%%%%%%%%%%%%%%%%%%%%%%%%%%%%%%%%%%%%%%%%%%%

The symbol

\[ \lor \]

is called the \emph{logical disjunction symbol} and is read

\begin{center}
``or.''
\end{center}

\begin{definitionbox}{Disjunction}
The disjunction
\[ P\lor Q \]
is true when at least one of \(P\) or \(Q\) is true.
\end{definitionbox}

In standard mathematical logic, ``or'' is inclusive.

Therefore \(P\lor Q\) is also true when both \(P\) and \(Q\) are true.

%%%%%%%%%%%%%%%%%%%%%%%%%%%%%%%%%%%%%%%%%%%%%%%%%%%%%%%%%%%%
\subsection{Truth Table for Disjunction}
\label{subsec:disjunction-truth-table}
%%%%%%%%%%%%%%%%%%%%%%%%%%%%%%%%%%%%%%%%%%%%%%%%%%%%%%%%%%%%

\[
\begin{array}{c|c|c}
P & Q & P\lor Q\\ \hline
T & T & T\\
T & F & T\\
F & T & T\\
F & F & F
\end{array}
\]

%%%%%%%%%%%%%%%%%%%%%%%%%%%%%%%%%%%%%%%%%%%%%%%%%%%%%%%%%%%%
\subsection{Inclusive versus Exclusive Or}
\label{subsec:inclusive-exclusive-or}
%%%%%%%%%%%%%%%%%%%%%%%%%%%%%%%%%%%%%%%%%%%%%%%%%%%%%%%%%%%%

Ordinary conversation sometimes uses ``or'' exclusively.

For example,

\begin{center}
``You may choose soup or salad.''
\end{center}

may imply that exactly one should be chosen.

Mathematical logic normally uses inclusive or:

\[ P\lor Q \]

means that one or both propositions are true.

When exactly one should be true, this is called \emph{exclusive or}.

%%%%%%%%%%%%%%%%%%%%%%%%%%%%%%%%%%%%%%%%%%%%%%%%%%%%%%%%%%%%
\subsection{Implication}
\label{subsec:implication}
%%%%%%%%%%%%%%%%%%%%%%%%%%%%%%%%%%%%%%%%%%%%%%%%%%%%%%%%%%%%

The symbol

\[ \to \]

is called the \emph{implication arrow} in logic.

The statement

\[ P\to Q \]

is read

\begin{center}
``if \(P\), then \(Q\)''
\end{center}

or

\begin{center}
``\(P\) implies \(Q\).''
\end{center}

\begin{definitionbox}{Implication}
The implication
\[ P\to Q \]
is false only when \(P\) is true and \(Q\) is false.
\end{definitionbox}

Here \(P\) is called the \emph{hypothesis} or \emph{antecedent}, and \(Q\)
is called the \emph{conclusion} or \emph{consequent}.

%%%%%%%%%%%%%%%%%%%%%%%%%%%%%%%%%%%%%%%%%%%%%%%%%%%%%%%%%%%%
\subsection{Truth Table for Implication}
\label{subsec:implication-truth-table}
%%%%%%%%%%%%%%%%%%%%%%%%%%%%%%%%%%%%%%%%%%%%%%%%%%%%%%%%%%%%

\[
\begin{array}{c|c|c}
P & Q & P\to Q\\ \hline
T & T & T\\
T & F & F\\
F & T & T\\
F & F & T
\end{array}
\]

The only false case is

\[ T\to F. \]

%%%%%%%%%%%%%%%%%%%%%%%%%%%%%%%%%%%%%%%%%%%%%%%%%%%%%%%%%%%%
\subsection{Why a False Hypothesis Gives a True Implication}
\label{subsec:false-hypothesis-implication}
%%%%%%%%%%%%%%%%%%%%%%%%%%%%%%%%%%%%%%%%%%%%%%%%%%%%%%%%%%%%

The truth table for implication sometimes feels unintuitive at first.

Consider the statement

\begin{center}
``If an integer is divisible by \(4\), then it is even.''
\end{center}

An integer such as \(6\) is not divisible by \(4\).

The statement does not claim anything about what must happen when the integer
is not divisible by \(4\).

Therefore \(6\) does not provide a counterexample.

To make an implication false, one must find a case where the hypothesis is
true and the conclusion is false.

%%%%%%%%%%%%%%%%%%%%%%%%%%%%%%%%%%%%%%%%%%%%%%%%%%%%%%%%%%%%
\subsection{Necessary and Sufficient Conditions}
\label{subsec:necessary-sufficient}
%%%%%%%%%%%%%%%%%%%%%%%%%%%%%%%%%%%%%%%%%%%%%%%%%%%%%%%%%%%%

The implication

\[ P\to Q \]

can be expressed in several equivalent ways:

\begin{itemize}
    \item if \(P\), then \(Q\);
    \item \(P\) implies \(Q\);
    \item \(P\) is sufficient for \(Q\);
    \item \(Q\) is necessary for \(P\);
    \item \(P\) only if \(Q\).
\end{itemize}

The distinction between ``if'' and ``only if'' is especially important.

%%%%%%%%%%%%%%%%%%%%%%%%%%%%%%%%%%%%%%%%%%%%%%%%%%%%%%%%%%%%
\subsection{Understanding ``Only If''}
\label{subsec:only-if}
%%%%%%%%%%%%%%%%%%%%%%%%%%%%%%%%%%%%%%%%%%%%%%%%%%%%%%%%%%%%

The statement

\begin{center}
``\(P\) only if \(Q\)''
\end{center}

means

\[ P\to Q. \]

Thus \(Q\) is required whenever \(P\) occurs.

For example,

\begin{center}
``An integer is divisible by \(4\) only if it is even''
\end{center}

means

\[ 4\mid n\to 2\mid n. \]

%%%%%%%%%%%%%%%%%%%%%%%%%%%%%%%%%%%%%%%%%%%%%%%%%%%%%%%%%%%%
\subsection{Converse}
\label{subsec:converse}
%%%%%%%%%%%%%%%%%%%%%%%%%%%%%%%%%%%%%%%%%%%%%%%%%%%%%%%%%%%%

Given

\[ P\to Q, \]

the \emph{converse} is

\[ Q\to P. \]

For example,

\begin{center}
Original: If a number is divisible by \(4\), then it is even.
\end{center}

The converse is

\begin{center}
If a number is even, then it is divisible by \(4\).
\end{center}

The original statement is true, but its converse is false.

For instance, \(6\) is even but not divisible by \(4\).

%%%%%%%%%%%%%%%%%%%%%%%%%%%%%%%%%%%%%%%%%%%%%%%%%%%%%%%%%%%%
\subsection{Inverse}
\label{subsec:logical-inverse}
%%%%%%%%%%%%%%%%%%%%%%%%%%%%%%%%%%%%%%%%%%%%%%%%%%%%%%%%%%%%

The \emph{inverse} of

\[ P\to Q \]

is

\[ \neg P\to\neg Q. \]

For the previous example:

\begin{center}
If a number is not divisible by \(4\), then it is not even.
\end{center}

This statement is false because \(6\) is not divisible by \(4\) but is even.

%%%%%%%%%%%%%%%%%%%%%%%%%%%%%%%%%%%%%%%%%%%%%%%%%%%%%%%%%%%%
\subsection{Contrapositive}
\label{subsec:contrapositive}
%%%%%%%%%%%%%%%%%%%%%%%%%%%%%%%%%%%%%%%%%%%%%%%%%%%%%%%%%%%%

The \emph{contrapositive} of

\[ P\to Q \]

is

\[ \neg Q\to\neg P. \]

For example,

\begin{center}
If an integer is not even, then it is not divisible by \(4\).
\end{center}

This statement is logically equivalent to the original implication.

%%%%%%%%%%%%%%%%%%%%%%%%%%%%%%%%%%%%%%%%%%%%%%%%%%%%%%%%%%%%
\subsection{Implication and Contrapositive}
\label{subsec:implication-contrapositive-equivalence}
%%%%%%%%%%%%%%%%%%%%%%%%%%%%%%%%%%%%%%%%%%%%%%%%%%%%%%%%%%%%

A fundamental logical equivalence is

\[ P\to Q\equiv\neg Q\to\neg P. \]

The symbol

\[ \equiv \]

is called the \emph{logical equivalence symbol} in this context and is read

\begin{center}
``is logically equivalent to.''
\end{center}

This equivalence is the foundation of proof by contraposition.

%%%%%%%%%%%%%%%%%%%%%%%%%%%%%%%%%%%%%%%%%%%%%%%%%%%%%%%%%%%%
\subsection{Converse and Inverse}
\label{subsec:converse-inverse-equivalence}
%%%%%%%%%%%%%%%%%%%%%%%%%%%%%%%%%%%%%%%%%%%%%%%%%%%%%%%%%%%%

The converse and inverse are logically equivalent to each other:

\[ Q\to P\equiv\neg P\to\neg Q. \]

However, neither is generally equivalent to the original implication.

%%%%%%%%%%%%%%%%%%%%%%%%%%%%%%%%%%%%%%%%%%%%%%%%%%%%%%%%%%%%
\subsection{Biconditional}
\label{subsec:biconditional}
%%%%%%%%%%%%%%%%%%%%%%%%%%%%%%%%%%%%%%%%%%%%%%%%%%%%%%%%%%%%

The symbol

\[ \leftrightarrow \]

is called the \emph{biconditional symbol} and is read

\begin{center}
``if and only if.''
\end{center}

\begin{definitionbox}{Biconditional}
The biconditional
\[ P\leftrightarrow Q \]
is true when \(P\) and \(Q\) have the same truth value.
\end{definitionbox}

It means both directions hold:

\[ P\to Q \]

and

\[ Q\to P. \]

%%%%%%%%%%%%%%%%%%%%%%%%%%%%%%%%%%%%%%%%%%%%%%%%%%%%%%%%%%%%
\subsection{Truth Table for Biconditional}
\label{subsec:biconditional-truth-table}
%%%%%%%%%%%%%%%%%%%%%%%%%%%%%%%%%%%%%%%%%%%%%%%%%%%%%%%%%%%%

\[
\begin{array}{c|c|c}
P & Q & P\leftrightarrow Q\\ \hline
T & T & T\\
T & F & F\\
F & T & F\\
F & F & T
\end{array}
\]

%%%%%%%%%%%%%%%%%%%%%%%%%%%%%%%%%%%%%%%%%%%%%%%%%%%%%%%%%%%%
\subsection{The Phrase ``If and Only If''}
\label{subsec:iff-language}
%%%%%%%%%%%%%%%%%%%%%%%%%%%%%%%%%%%%%%%%%%%%%%%%%%%%%%%%%%%%

The phrase

\begin{center}
``if and only if''
\end{center}

is often abbreviated

\begin{center}
``iff.''
\end{center}

Thus,

\[ P\leftrightarrow Q \]

may be read

\begin{center}
``\(P\) iff \(Q\).''
\end{center}

A biconditional establishes both necessity and sufficiency.

%%%%%%%%%%%%%%%%%%%%%%%%%%%%%%%%%%%%%%%%%%%%%%%%%%%%%%%%%%%%
\subsection{Compound Propositions}
\label{subsec:compound-propositions}
%%%%%%%%%%%%%%%%%%%%%%%%%%%%%%%%%%%%%%%%%%%%%%%%%%%%%%%%%%%%

Logical connectives may be combined to create more complicated statements.

For example,

\[ (P\land Q)\to R. \]

Parentheses specify how the statement is grouped.

As in arithmetic, changing the grouping may change the meaning.

%%%%%%%%%%%%%%%%%%%%%%%%%%%%%%%%%%%%%%%%%%%%%%%%%%%%%%%%%%%%
\subsection{Constructing Truth Tables}
\label{subsec:constructing-truth-tables}
%%%%%%%%%%%%%%%%%%%%%%%%%%%%%%%%%%%%%%%%%%%%%%%%%%%%%%%%%%%%

A truth table for two independent propositions requires

\[ 2^2=4 \]

rows.

For three propositions, it requires

\[ 2^3=8 \]

rows.

More generally, \(n\) independent proposition variables require

\[ 2^n \]

possible combinations of truth values.

%%%%%%%%%%%%%%%%%%%%%%%%%%%%%%%%%%%%%%%%%%%%%%%%%%%%%%%%%%%%
\subsection{Worked Example: Compound Truth Table}
\label{subsec:worked-compound-truth-table}
%%%%%%%%%%%%%%%%%%%%%%%%%%%%%%%%%%%%%%%%%%%%%%%%%%%%%%%%%%%%

\begin{workedexamplebox}{Construct a Truth Table for \((P\land Q)\to P\)}

\[
\begin{array}{c|c|c|c}
P & Q & P\land Q & (P\land Q)\to P\\ \hline
T & T & T & T\\
T & F & F & T\\
F & T & F & T\\
F & F & F & T
\end{array}
\]

\begin{answerbox}
The statement is true for every possible truth assignment.
\end{answerbox}
\end{workedexamplebox}

%%%%%%%%%%%%%%%%%%%%%%%%%%%%%%%%%%%%%%%%%%%%%%%%%%%%%%%%%%%%
\subsection{Tautologies}
\label{subsec:tautologies}
%%%%%%%%%%%%%%%%%%%%%%%%%%%%%%%%%%%%%%%%%%%%%%%%%%%%%%%%%%%%

\begin{definitionbox}{Tautology}
A \emph{tautology} is a compound proposition that is true for every possible
assignment of truth values.
\end{definitionbox}

For example,

\[ P\lor\neg P \]

is always true.

This is called the \emph{law of excluded middle} in classical logic.

%%%%%%%%%%%%%%%%%%%%%%%%%%%%%%%%%%%%%%%%%%%%%%%%%%%%%%%%%%%%
\subsection{Contradictions}
\label{subsec:contradictions}
%%%%%%%%%%%%%%%%%%%%%%%%%%%%%%%%%%%%%%%%%%%%%%%%%%%%%%%%%%%%

\begin{definitionbox}{Contradiction}
A \emph{contradiction} is a compound proposition that is false for every
possible assignment of truth values.
\end{definitionbox}

For example,

\[ P\land\neg P \]

is always false.

%%%%%%%%%%%%%%%%%%%%%%%%%%%%%%%%%%%%%%%%%%%%%%%%%%%%%%%%%%%%
\subsection{Contingencies}
\label{subsec:contingencies}
%%%%%%%%%%%%%%%%%%%%%%%%%%%%%%%%%%%%%%%%%%%%%%%%%%%%%%%%%%%%

\begin{definitionbox}{Contingency}
A \emph{contingency} is a proposition that is true for some truth assignments
and false for others.
\end{definitionbox}

For example,

\[ P\land Q \]

is a contingency.

%%%%%%%%%%%%%%%%%%%%%%%%%%%%%%%%%%%%%%%%%%%%%%%%%%%%%%%%%%%%
\subsection{Logical Equivalence}
\label{subsec:logical-equivalence}
%%%%%%%%%%%%%%%%%%%%%%%%%%%%%%%%%%%%%%%%%%%%%%%%%%%%%%%%%%%%

Two propositions are logically equivalent if they have the same truth value
under every possible truth assignment.

We write

\[ P\equiv Q. \]

For example,

\[ P\to Q\equiv\neg P\lor Q. \]

This identity is useful because an implication can be rewritten using only
negation and disjunction.

%%%%%%%%%%%%%%%%%%%%%%%%%%%%%%%%%%%%%%%%%%%%%%%%%%%%%%%%%%%%
\subsection{Double Negation}
\label{subsec:double-negation}
%%%%%%%%%%%%%%%%%%%%%%%%%%%%%%%%%%%%%%%%%%%%%%%%%%%%%%%%%%%%

A proposition and its double negation are logically equivalent:

\[ \neg(\neg P)\equiv P. \]

Thus ``not not \(P\)'' has the same truth value as \(P\).

%%%%%%%%%%%%%%%%%%%%%%%%%%%%%%%%%%%%%%%%%%%%%%%%%%%%%%%%%%%%
\subsection{De Morgan's Laws for Logic}
\label{subsec:de-morgan-logic}
%%%%%%%%%%%%%%%%%%%%%%%%%%%%%%%%%%%%%%%%%%%%%%%%%%%%%%%%%%%%

Two fundamental logical identities are

\[ \neg(P\land Q)\equiv\neg P\lor\neg Q, \]

and

\[ \neg(P\lor Q)\equiv\neg P\land\neg Q. \]

These are the logical versions of the set identities introduced in
\cref{subsec:set-de-morgan-laws}.

%%%%%%%%%%%%%%%%%%%%%%%%%%%%%%%%%%%%%%%%%%%%%%%%%%%%%%%%%%%%
\subsection{Negating Everyday Mathematical Statements}
\label{subsec:negating-statements}
%%%%%%%%%%%%%%%%%%%%%%%%%%%%%%%%%%%%%%%%%%%%%%%%%%%%%%%%%%%%

Consider

\begin{center}
``\(x>2\) and \(x<7\).''
\end{center}

Its negation is

\begin{center}
``\(x\leq2\) or \(x\geq7\).''
\end{center}

Symbolically,

\[ \neg(x>2\land x<7)\equiv x\leq2\lor x\geq7. \]

This follows from De Morgan's laws.

%%%%%%%%%%%%%%%%%%%%%%%%%%%%%%%%%%%%%%%%%%%%%%%%%%%%%%%%%%%%
\subsection{Predicates}
\label{subsec:predicates}
%%%%%%%%%%%%%%%%%%%%%%%%%%%%%%%%%%%%%%%%%%%%%%%%%%%%%%%%%%%%

\begin{definitionbox}{Predicate}
A \emph{predicate} is a statement containing one or more variables whose
truth value depends on the values assigned to those variables.
\end{definitionbox}

For example,

\[ P(x):x^2>4 \]

is a predicate on \(\R\).

It is true for some values of \(x\) and false for others.

For instance,

\[ P(3) \]

is true, while

\[ P(1) \]

is false.

%%%%%%%%%%%%%%%%%%%%%%%%%%%%%%%%%%%%%%%%%%%%%%%%%%%%%%%%%%%%
\subsection{Universal Quantifier}
\label{subsec:universal-quantifier}
%%%%%%%%%%%%%%%%%%%%%%%%%%%%%%%%%%%%%%%%%%%%%%%%%%%%%%%%%%%%

The symbol

\[ \forall \]

is called the \emph{universal quantifier} and is read

\begin{center}
``for all''
\end{center}

or

\begin{center}
``for every.''
\end{center}

For example,

\[ \forall x\in\R,\;x^2\geq0 \]

is read

\begin{center}
``For every real number \(x\), \(x^2\) is greater than or equal to zero.''
\end{center}

%%%%%%%%%%%%%%%%%%%%%%%%%%%%%%%%%%%%%%%%%%%%%%%%%%%%%%%%%%%%
\subsection{Existential Quantifier}
\label{subsec:existential-quantifier}
%%%%%%%%%%%%%%%%%%%%%%%%%%%%%%%%%%%%%%%%%%%%%%%%%%%%%%%%%%%%

The symbol

\[ \exists \]

is called the \emph{existential quantifier} and is read

\begin{center}
``there exists.''
\end{center}

For example,

\[ \exists x\in\Z\text{ such that }x^2=9 \]

is read

\begin{center}
``There exists an integer \(x\) such that \(x^2=9\).''
\end{center}

This statement is true because \(x=3\) and \(x=-3\) both work.

%%%%%%%%%%%%%%%%%%%%%%%%%%%%%%%%%%%%%%%%%%%%%%%%%%%%%%%%%%%%
\subsection{Unique Existence}
\label{subsec:unique-existence}
%%%%%%%%%%%%%%%%%%%%%%%%%%%%%%%%%%%%%%%%%%%%%%%%%%%%%%%%%%%%

The symbol

\[ \exists! \]

is read

\begin{center}
``there exists exactly one.''
\end{center}

For example,

\[ \exists!x\in\R\text{ such that }x+4=7. \]

This statement is true because the unique solution is

\[ x=3. \]

%%%%%%%%%%%%%%%%%%%%%%%%%%%%%%%%%%%%%%%%%%%%%%%%%%%%%%%%%%%%
\subsection{Quantifier Domains Matter}
\label{subsec:quantifier-domains}
%%%%%%%%%%%%%%%%%%%%%%%%%%%%%%%%%%%%%%%%%%%%%%%%%%%%%%%%%%%%

The truth of a quantified statement may depend on its domain.

For example,

\[ \exists x\in\R\text{ such that }x^2=2 \]

is true.

But

\[ \exists x\in\Q\text{ such that }x^2=2 \]

is false.

Thus the universe over which a variable ranges is part of the statement.

%%%%%%%%%%%%%%%%%%%%%%%%%%%%%%%%%%%%%%%%%%%%%%%%%%%%%%%%%%%%
\subsection{Negating Universal Statements}
\label{subsec:negating-universal}
%%%%%%%%%%%%%%%%%%%%%%%%%%%%%%%%%%%%%%%%%%%%%%%%%%%%%%%%%%%%

A universal statement is negated by changing

\[ \forall \]

to

\[ \exists \]

and negating the predicate.

Thus,

\[ \neg(\forall x,\;P(x))\equiv\exists x,\;\neg P(x). \]

In words:

\begin{center}
``It is not true that every \(x\) satisfies \(P\)''
\end{center}

means

\begin{center}
``There exists at least one \(x\) that does not satisfy \(P\).''
\end{center}

%%%%%%%%%%%%%%%%%%%%%%%%%%%%%%%%%%%%%%%%%%%%%%%%%%%%%%%%%%%%
\subsection{Negating Existential Statements}
\label{subsec:negating-existential}
%%%%%%%%%%%%%%%%%%%%%%%%%%%%%%%%%%%%%%%%%%%%%%%%%%%%%%%%%%%%

Similarly,

\[ \neg(\exists x,\;P(x))\equiv\forall x,\;\neg P(x). \]

Thus,

\begin{center}
``There does not exist an \(x\) satisfying \(P\)''
\end{center}

means

\begin{center}
``Every \(x\) fails to satisfy \(P\).''
\end{center}

%%%%%%%%%%%%%%%%%%%%%%%%%%%%%%%%%%%%%%%%%%%%%%%%%%%%%%%%%%%%
\subsection{Worked Example: Negating a Quantified Statement}
\label{subsec:worked-negating-quantifier}
%%%%%%%%%%%%%%%%%%%%%%%%%%%%%%%%%%%%%%%%%%%%%%%%%%%%%%%%%%%%

\begin{workedexamplebox}{Negate a Universal Statement}

Negate

\[ \forall x\in\R,\;x^2+x+1>0. \]

\begin{tabularx}{\textwidth}{@{}>{\raggedright\arraybackslash}p{0.42\textwidth}X@{}}
Change the universal quantifier. & \(\forall\to\exists\)\\
Negate the inequality. & \(x^2+x+1\leq0\)\\
Keep the same domain. & \(x\in\R\)
\end{tabularx}

\begin{answerbox}
\[ \exists x\in\R\text{ such that }x^2+x+1\leq0 \]
\end{answerbox}
\end{workedexamplebox}

%%%%%%%%%%%%%%%%%%%%%%%%%%%%%%%%%%%%%%%%%%%%%%%%%%%%%%%%%%%%
\subsection{Multiple Quantifiers}
\label{subsec:multiple-quantifiers}
%%%%%%%%%%%%%%%%%%%%%%%%%%%%%%%%%%%%%%%%%%%%%%%%%%%%%%%%%%%%

Statements may contain several quantifiers.

For example,

\[ \forall x\in\R,\;\exists y\in\R\text{ such that }y>x. \]

This is read

\begin{center}
``For every real number \(x\), there exists a real number \(y\) such that
\(y>x\).''
\end{center}

The statement is true.

For any \(x\), one may choose

\[ y=x+1. \]

%%%%%%%%%%%%%%%%%%%%%%%%%%%%%%%%%%%%%%%%%%%%%%%%%%%%%%%%%%%%
\subsection{Order of Quantifiers Matters}
\label{subsec:quantifier-order}
%%%%%%%%%%%%%%%%%%%%%%%%%%%%%%%%%%%%%%%%%%%%%%%%%%%%%%%%%%%%

Consider

\[ \forall x\in\R,\;\exists y\in\R\text{ such that }y>x. \]

This is true.

Now reverse the quantifiers:

\[ \exists y\in\R,\;\forall x\in\R,\;y>x. \]

This would say that there is one real number larger than every real number.

That statement is false.

Therefore the order of quantifiers can fundamentally change meaning.

%%%%%%%%%%%%%%%%%%%%%%%%%%%%%%%%%%%%%%%%%%%%%%%%%%%%%%%%%%%%
\subsection{Translating Quantified Statements}
\label{subsec:translating-quantified-statements}
%%%%%%%%%%%%%%%%%%%%%%%%%%%%%%%%%%%%%%%%%%%%%%%%%%%%%%%%%%%%

The statement

\[ \forall x\in A,\;P(x) \]

means that every element of \(A\) satisfies \(P\).

The statement

\[ \exists x\in A,\;P(x) \]

means that at least one element of \(A\) satisfies \(P\).

The statement

\[ \exists!x\in A,\;P(x) \]

means exactly one element of \(A\) satisfies \(P\).

%%%%%%%%%%%%%%%%%%%%%%%%%%%%%%%%%%%%%%%%%%%%%%%%%%%%%%%%%%%%
\subsection{Arguments}
\label{subsec:logical-arguments}
%%%%%%%%%%%%%%%%%%%%%%%%%%%%%%%%%%%%%%%%%%%%%%%%%%%%%%%%%%%%

An \emph{argument} consists of propositions called premises together with a
conclusion.

A valid argument is one in which the conclusion must be true whenever all
premises are true.

Validity concerns logical form.

It does not require the premises themselves to be true in the real world.

%%%%%%%%%%%%%%%%%%%%%%%%%%%%%%%%%%%%%%%%%%%%%%%%%%%%%%%%%%%%
\subsection{Therefore Symbol}
\label{subsec:therefore-symbol}
%%%%%%%%%%%%%%%%%%%%%%%%%%%%%%%%%%%%%%%%%%%%%%%%%%%%%%%%%%%%

The symbol

\[ \therefore \]

is called the \emph{therefore symbol} and is read

\begin{center}
``therefore.''
\end{center}

It is sometimes used to introduce the conclusion of an argument.

For example,

\[ P,\qquad P\to Q,\qquad \therefore Q. \]

%%%%%%%%%%%%%%%%%%%%%%%%%%%%%%%%%%%%%%%%%%%%%%%%%%%%%%%%%%%%
\subsection{Modus Ponens}
\label{subsec:modus-ponens}
%%%%%%%%%%%%%%%%%%%%%%%%%%%%%%%%%%%%%%%%%%%%%%%%%%%%%%%%%%%%

\begin{definitionbox}{Modus Ponens}
The argument form
\[ P,\qquad P\to Q,\qquad \therefore Q \]
is called \emph{modus ponens}.
\end{definitionbox}

It may be read:

\begin{center}
\(P\) is true. If \(P\), then \(Q\). Therefore \(Q\).
\end{center}

This is one of the most basic valid forms of mathematical reasoning.

%%%%%%%%%%%%%%%%%%%%%%%%%%%%%%%%%%%%%%%%%%%%%%%%%%%%%%%%%%%%
\subsection{Modus Tollens}
\label{subsec:modus-tollens}
%%%%%%%%%%%%%%%%%%%%%%%%%%%%%%%%%%%%%%%%%%%%%%%%%%%%%%%%%%%%

\begin{definitionbox}{Modus Tollens}
The argument form
\[ P\to Q,\qquad \neg Q,\qquad \therefore\neg P \]
is called \emph{modus tollens}.
\end{definitionbox}

It relies on the equivalence between an implication and its contrapositive.

%%%%%%%%%%%%%%%%%%%%%%%%%%%%%%%%%%%%%%%%%%%%%%%%%%%%%%%%%%%%
\subsection{Hypothetical Syllogism}
\label{subsec:hypothetical-syllogism}
%%%%%%%%%%%%%%%%%%%%%%%%%%%%%%%%%%%%%%%%%%%%%%%%%%%%%%%%%%%%

Another valid argument form is

\[ P\to Q,\qquad Q\to R,\qquad \therefore P\to R. \]

This is called a \emph{hypothetical syllogism}.

It expresses transitivity of implication.

%%%%%%%%%%%%%%%%%%%%%%%%%%%%%%%%%%%%%%%%%%%%%%%%%%%%%%%%%%%%
\subsection{Disjunctive Syllogism}
\label{subsec:disjunctive-syllogism}
%%%%%%%%%%%%%%%%%%%%%%%%%%%%%%%%%%%%%%%%%%%%%%%%%%%%%%%%%%%%

The argument

\[ P\lor Q,\qquad \neg P,\qquad \therefore Q \]

is called a \emph{disjunctive syllogism}.

If at least one of \(P\) or \(Q\) is true and \(P\) is false, then \(Q\)
must be true.

%%%%%%%%%%%%%%%%%%%%%%%%%%%%%%%%%%%%%%%%%%%%%%%%%%%%%%%%%%%%
\subsection{Invalid Form: Affirming the Consequent}
\label{subsec:affirming-consequent}
%%%%%%%%%%%%%%%%%%%%%%%%%%%%%%%%%%%%%%%%%%%%%%%%%%%%%%%%%%%%

The argument form

\[ P\to Q,\qquad Q,\qquad \therefore P \]

is invalid.

This error is called \emph{affirming the consequent}.

For example,

\begin{center}
If a number is divisible by \(4\), then it is even.

\(6\) is even.

Therefore \(6\) is divisible by \(4\).
\end{center}

The conclusion is false.

%%%%%%%%%%%%%%%%%%%%%%%%%%%%%%%%%%%%%%%%%%%%%%%%%%%%%%%%%%%%
\subsection{Invalid Form: Denying the Antecedent}
\label{subsec:denying-antecedent}
%%%%%%%%%%%%%%%%%%%%%%%%%%%%%%%%%%%%%%%%%%%%%%%%%%%%%%%%%%%%

The argument form

\[ P\to Q,\qquad \neg P,\qquad \therefore\neg Q \]

is also invalid.

This is called \emph{denying the antecedent}.

For example,

\begin{center}
If an integer is divisible by \(4\), then it is even.

\(6\) is not divisible by \(4\).

Therefore \(6\) is not even.
\end{center}

Again, the conclusion is false.

%%%%%%%%%%%%%%%%%%%%%%%%%%%%%%%%%%%%%%%%%%%%%%%%%%%%%%%%%%%%
\subsection{Counterexamples}
\label{subsec:logic-counterexamples}
%%%%%%%%%%%%%%%%%%%%%%%%%%%%%%%%%%%%%%%%%%%%%%%%%%%%%%%%%%%%

A universal statement has the form

\[ \forall x\in A,\;P(x). \]

To show that it is false, it is enough to find one element \(a\in A\) such
that

\[ \neg P(a). \]

Such an example is called a \emph{counterexample}.

For instance, the statement

\begin{center}
``Every prime number is odd''
\end{center}

is false because

\[ 2 \]

is prime and even.

%%%%%%%%%%%%%%%%%%%%%%%%%%%%%%%%%%%%%%%%%%%%%%%%%%%%%%%%%%%%
\subsection{Logic and Sets}
\label{subsec:logic-and-sets}
%%%%%%%%%%%%%%%%%%%%%%%%%%%%%%%%%%%%%%%%%%%%%%%%%%%%%%%%%%%%

Logical operations closely parallel set operations.

If \(P(x)\) and \(Q(x)\) are predicates, define

\[ A=\{x:P(x)\},\qquad B=\{x:Q(x)\}. \]

Then

\[ P(x)\lor Q(x) \]

corresponds to

\[ A\cup B. \]

Similarly,

\[ P(x)\land Q(x) \]

corresponds to

\[ A\cap B. \]

Negation corresponds to set complement.

%%%%%%%%%%%%%%%%%%%%%%%%%%%%%%%%%%%%%%%%%%%%%%%%%%%%%%%%%%%%
\subsection{Logic and De Morgan's Laws}
\label{subsec:logic-set-de-morgan}
%%%%%%%%%%%%%%%%%%%%%%%%%%%%%%%%%%%%%%%%%%%%%%%%%%%%%%%%%%%%

The logical identity

\[ \neg(P\lor Q)\equiv\neg P\land\neg Q \]

corresponds to

\[ (A\cup B)^c=A^c\cap B^c. \]

Likewise,

\[ \neg(P\land Q)\equiv\neg P\lor\neg Q \]

corresponds to

\[ (A\cap B)^c=A^c\cup B^c. \]

This connection is one reason logical reasoning is so useful in set proofs.

%%%%%%%%%%%%%%%%%%%%%%%%%%%%%%%%%%%%%%%%%%%%%%%%%%%%%%%%%%%%
\subsection{Logic and Proof}
\label{subsec:logic-and-proof}
%%%%%%%%%%%%%%%%%%%%%%%%%%%%%%%%%%%%%%%%%%%%%%%%%%%%%%%%%%%%

Different logical forms suggest different proof strategies.

To prove

\[ P\to Q, \]

one may use a direct proof or prove the contrapositive.

To disprove

\[ \forall x,\;P(x), \]

one may provide a counterexample.

To prove

\[ P\leftrightarrow Q, \]

one usually proves both

\[ P\to Q \]

and

\[ Q\to P. \]

These methods are developed systematically in
\cref{sec:methods-of-proof}.

%%%%%%%%%%%%%%%%%%%%%%%%%%%%%%%%%%%%%%%%%%%%%%%%%%%%%%%%%%%%
\subsection{Common Logic Errors}
\label{subsec:common-logic-errors}
%%%%%%%%%%%%%%%%%%%%%%%%%%%%%%%%%%%%%%%%%%%%%%%%%%%%%%%%%%%%

\subsubsection{Assuming the Converse Is Automatically True}

From

\[ P\to Q \]

one may not conclude

\[ Q\to P. \]

The converse requires a separate argument.

%%%%%%%%%%%%%%%%%%%%%%%%%%%%%%%%%%%%%%%%%%%%%%%%%%%%%%%%%%%%
\subsubsection{Confusing Necessary and Sufficient}
\label{subsec:error-necessary-sufficient}
%%%%%%%%%%%%%%%%%%%%%%%%%%%%%%%%%%%%%%%%%%%%%%%%%%%%%%%%%%%%

If

\[ P\to Q, \]

then \(P\) is sufficient for \(Q\), while \(Q\) is necessary for \(P\).

Reversing these statements changes the logical meaning.

%%%%%%%%%%%%%%%%%%%%%%%%%%%%%%%%%%%%%%%%%%%%%%%%%%%%%%%%%%%%
\subsubsection{Negating Only Part of a Statement}
\label{subsec:error-partial-negation}
%%%%%%%%%%%%%%%%%%%%%%%%%%%%%%%%%%%%%%%%%%%%%%%%%%%%%%%%%%%%

The negation of

\[ P\land Q \]

is not

\[ \neg P\land\neg Q. \]

The correct negation is

\[ \neg P\lor\neg Q. \]

%%%%%%%%%%%%%%%%%%%%%%%%%%%%%%%%%%%%%%%%%%%%%%%%%%%%%%%%%%%%
\subsubsection{Forgetting to Change the Quantifier}
\label{subsec:error-negating-quantifier}
%%%%%%%%%%%%%%%%%%%%%%%%%%%%%%%%%%%%%%%%%%%%%%%%%%%%%%%%%%%%

The negation of

\[ \forall x,\;P(x) \]

is

\[ \exists x,\;\neg P(x), \]

not

\[ \forall x,\;\neg P(x). \]

%%%%%%%%%%%%%%%%%%%%%%%%%%%%%%%%%%%%%%%%%%%%%%%%%%%%%%%%%%%%
\subsubsection{Ignoring Quantifier Order}
\label{subsec:error-quantifier-order}
%%%%%%%%%%%%%%%%%%%%%%%%%%%%%%%%%%%%%%%%%%%%%%%%%%%%%%%%%%%%

The statements

\[ \forall x\exists y\,P(x,y) \]

and

\[ \exists y\forall x\,P(x,y) \]

usually have different meanings.

Their quantifiers cannot generally be exchanged.

%%%%%%%%%%%%%%%%%%%%%%%%%%%%%%%%%%%%%%%%%%%%%%%%%%%%%%%%%%%%
\subsubsection{Treating Evidence as Proof}
\label{subsec:error-evidence-proof}
%%%%%%%%%%%%%%%%%%%%%%%%%%%%%%%%%%%%%%%%%%%%%%%%%%%%%%%%%%%%

Checking many examples can suggest that a universal claim is true, but examples
alone do not prove it.

A proof must establish that the claim holds in every required case.

%%%%%%%%%%%%%%%%%%%%%%%%%%%%%%%%%%%%%%%%%%%%%%%%%%%%%%%%%%%%
\subsection{Worked Examples}
\label{subsec:logic-worked-examples}
%%%%%%%%%%%%%%%%%%%%%%%%%%%%%%%%%%%%%%%%%%%%%%%%%%%%%%%%%%%%

\begin{workedexamplebox}{Write a Contrapositive}

Consider the statement:

\begin{center}
If \(n^2\) is odd, then \(n\) is odd.
\end{center}

\begin{tabularx}{\textwidth}{@{}>{\raggedright\arraybackslash}p{0.42\textwidth}X@{}}
Identify \(P\). & \(n^2\text{ is odd}\)\\
Identify \(Q\). & \(n\text{ is odd}\)\\
Negate \(Q\). & \(n\text{ is even}\)\\
Negate \(P\). & \(n^2\text{ is even}\)\\
Reverse the implication. & \(\neg Q\to\neg P\)
\end{tabularx}

\begin{answerbox}
If \(n\) is even, then \(n^2\) is even.
\end{answerbox}
\end{workedexamplebox}

\begin{workedexamplebox}{Negate a Statement with Two Quantifiers}

Negate

\[ \forall x\in\R,\;\exists y\in\R\text{ such that }y>x. \]

\begin{tabularx}{\textwidth}{@{}>{\raggedright\arraybackslash}p{0.42\textwidth}X@{}}
Negate the first quantifier. & \(\forall\to\exists\)\\
Negate the second quantifier. & \(\exists\to\forall\)\\
Negate \(y>x\). & \(y\leq x\)
\end{tabularx}

\begin{answerbox}
\[ \exists x\in\R,\;\forall y\in\R,\;y\leq x \]
\end{answerbox}
\end{workedexamplebox}

\begin{workedexamplebox}{Identify an Invalid Argument}

Consider

\[ P\to Q,\qquad Q,\qquad \therefore P. \]

\begin{tabularx}{\textwidth}{@{}>{\raggedright\arraybackslash}p{0.42\textwidth}X@{}}
The implication gives one direction only. & \(P\to Q\)\\
The conclusion reverses that direction. & \(Q\to P\) is being assumed\\
This form is not valid. & Affirming the consequent
\end{tabularx}

\begin{answerbox}
The argument is invalid.
\end{answerbox}
\end{workedexamplebox}

%%%%%%%%%%%%%%%%%%%%%%%%%%%%%%%%%%%%%%%%%%%%%%%%%%%%%%%%%%%%
\subsection{Exercises}
\label{subsec:logic-exercises}
%%%%%%%%%%%%%%%%%%%%%%%%%%%%%%%%%%%%%%%%%%%%%%%%%%%%%%%%%%%%

\begin{exercise}[1]
Determine whether each statement is a proposition:
\begin{enumerate}
    \item \(7+4=11\);
    \item \(x+3=8\);
    \item ``Close the door.'';
    \item \(5>9\).
\end{enumerate}
\end{exercise}

\begin{exercise}[1]
If \(P\) is true and \(Q\) is false, determine the truth value of:
\[ \neg P,\qquad P\land Q,\qquad P\lor Q,\qquad P\to Q,\qquad Q\to P. \]
\end{exercise}

\begin{exercise}[1]
Construct the truth table for
\[ \neg(P\land Q). \]
\end{exercise}

\begin{exercise}[1]
Construct the truth table for
\[ \neg P\lor Q. \]
\end{exercise}

\begin{exercise}[2]
Use a truth table to verify
\[ P\to Q\equiv\neg P\lor Q. \]
\end{exercise}

\begin{exercise}[2]
Use a truth table to verify
\[ P\to Q\equiv\neg Q\to\neg P. \]
\end{exercise}

\begin{exercise}[1]
Write the converse, inverse, and contrapositive of:

\begin{center}
If an integer is divisible by \(6\), then it is divisible by \(3\).
\end{center}
\end{exercise}

\begin{exercise}[1]
Translate
\[ P\to Q \]
using the words ``sufficient'' and ``necessary.''
\end{exercise}

\begin{exercise}[1]
Determine whether
\[ P\lor\neg P \]
is a tautology, contradiction, or contingency.
\end{exercise}

\begin{exercise}[1]
Determine whether
\[ P\land\neg P \]
is a tautology, contradiction, or contingency.
\end{exercise}

\begin{exercise}[2]
Negate:
\[ x>2\text{ and }x<8. \]
\end{exercise}

\begin{exercise}[2]
Negate:
\[ x\leq-1\text{ or }x\geq5. \]
\end{exercise}

\begin{exercise}[1]
Write in words:
\[ \forall x\in\R,\;x^2\geq0. \]
\end{exercise}

\begin{exercise}[1]
Write symbolically:

\begin{center}
There exists an integer whose square equals \(16\).
\end{center}
\end{exercise}

\begin{exercise}[2]
Negate:
\[ \forall n\in\Z,\;n^2\geq0. \]
\end{exercise}

\begin{exercise}[2]
Negate:
\[ \exists x\in\R\text{ such that }x^2+1=0. \]
\end{exercise}

\begin{exercise}[2]
Explain why
\[ \forall x\in\R,\;\exists y\in\R,\;y>x \]
is true.
\end{exercise}

\begin{exercise}[2]
Explain why
\[ \exists y\in\R,\;\forall x\in\R,\;y>x \]
is false.
\end{exercise}

\begin{exercise}[2]
Identify the argument form:
\[ P,\qquad P\to Q,\qquad \therefore Q. \]
\end{exercise}

\begin{exercise}[2]
Identify the error in:
\[ P\to Q,\qquad Q,\qquad \therefore P. \]
\end{exercise}

\begin{exercise}[2]
Give a counterexample to:

\begin{center}
Every positive integer is prime.
\end{center}
\end{exercise}

\begin{exercise}[3]
Prove using a truth table that
\[ \neg(P\lor Q)\equiv\neg P\land\neg Q. \]
\end{exercise}

\begin{exercise}[3]
Prove using a truth table that
\[ P\leftrightarrow Q\equiv(P\to Q)\land(Q\to P). \]
\end{exercise}

\begin{exercise}[3]
Negate the statement
\[ \forall x\in\R,\;\exists y\in\R,\;\forall z\in\R,\;P(x,y,z). \]
\end{exercise}

%%%%%%%%%%%%%%%%%%%%%%%%%%%%%%%%%%%%%%%%%%%%%%%%%%%%%%%%%%%%
\subsection{Section Connections}
\label{subsec:logic-connections}
%%%%%%%%%%%%%%%%%%%%%%%%%%%%%%%%%%%%%%%%%%%%%%%%%%%%%%%%%%%%

\chapterconnection{
Logic supplies the formal structure behind the proof methods developed in
\cref{sec:methods-of-proof}. Logical conjunction, disjunction, and negation
parallel intersection, union, and complement from
\cref{sec:elementary-set-theory}. Quantifiers occur throughout analysis,
algebra, topology, probability, and nearly every later chapter whenever
mathematical statements describe all objects or assert the existence of
particular objects.
}

%%%%%%%%%%%%%%%%%%%%%%%%%%%%%%%%%%%%%%%%%%%%%%%%%%%%%%%%%%%%
% SECTION SOURCE TRACKING
%%%%%%%%%%%%%%%%%%%%%%%%%%%%%%%%%%%%%%%%%%%%%%%%%%%%%%%%%%%%

%------------------------------------------------------------
% 0.7 Elementary Logic
%------------------------------------------------------------
%
% Sources:
%   - Richard Hammack, Book of Proof, Third Edition.
%     Key: hammack2018bookofproof
%
% Used for:
%   - propositions and logical connectives;
%   - implication, converse, inverse, and contrapositive;
%   - truth tables and logical equivalence;
%   - predicates and quantifiers;
%   - negation of quantified statements;
%   - valid and invalid argument forms;
%   - relationship between logic and proof.
%
% Notes:
%   - Formal symbolic logic is developed in greater depth in Chapter 12.
%   - Proof applications are developed in the following section.
%
\nocite{hammack2018bookofproof}